\chapter{Generative Bayesian Computation for Vertical Counting}
\label{ch:gbc}

Generative Bayesian computation, GBC, is the family of methods in which a posterior or other target distribution is represented not by a density but by a generator: a map from a base distribution, usually uniform or Gauss\index{Gauss, Carl Friedrich}ian noise, to the target, learned from simulated pairs by quantile regression or a related loss. It was developed by the authors and coauthors \cite{PolsonSokolovGBC} for likelihood free Bayesian inference, and its central object is the quantile function, which is also the central object of this book. This chapter shows that GBC gives a practical implementation of vertical counting, states the theorems that govern the implementation, and explains what the generator view adds to the analysis of the level set sampler\index{level set sampler}.

\section{Generators and quantile functions}

A generator for a distribution $Q$ on $\mathbb{R}^d$ is a measurable map $G$ from a base space with a fixed law $\mu$ to $\mathbb{R}^d$ with $G(\xi) \sim Q$ for $\xi \sim \mu$. In one dimension with $\mu$ uniform on $[0, 1]$ the canonical generator is the quantile function\index{quantile function} $F_Q^{-1}$, and it is the unique nondecreasing one. In dimension $d$ the Rosenblatt transform, which composes conditional quantile functions coordinate by coordinate, is a generator from the uniform distribution on $[0, 1]^d$, and the noise outsourcing lemma of probability theory says that any conditional law $Q(\cdot \mid y)$ can be written as $G(y, \xi)$ for a single map $G$ and a single uniform $\xi$.

The GBC method learns $G$ from samples. Given pairs $(y_i, \theta_i)$ drawn from the joint law of data and parameter, one fits a network $G_\phi(y, \xi)$ to minimise a loss whose expectation is minimised exactly when $G_\phi(y, \xi) \sim Q(\cdot \mid y)$: the pinball or check loss for quantile regression in one dimension, an energy or kernel score in several. The fitted $G_\phi$ is then a sampler for the posterior at any observed $y$ with no further simulation. The method is likelihood free, since only forward samples are needed, and it is the natural implementation of the vertical approach because what it learns is the quantile function.

\section{The vertical counting generator}

Apply this to counting. The quantity is $Z = \int_0^1 \Lambda(s) \, ds$ with $\Lambda$ the quantile function of $Y = L(X)$ under $P$. A generator for $Y$ from the uniform base is $\Lambda$ itself, up to the reversal $s \mapsto 1 - s$. So the GBC estimate of $Z$ is: draw $X_1, \dots, X_M$ from $P$, compute $Y_i = L(X_i)$, fit a monotone network $\hat\Lambda_\phi$ to the empirical quantiles of the $Y_i$ by the pinball loss, and integrate $\hat\Lambda_\phi$ over $[0, 1]$ by the stratified\index{stratified sampling} rule of Chapter~\ref{ch:beyond} or in closed form if the network is piecewise linear.

For the rare event case this is useless as stated, since the $Y_i$ are almost all zero and the empirical quantile function is zero everywhere except on a fraction $\rho$ of the interval that a sample of size $M \ll 1/\rho$ never sees. The remedy is the ladder: fit the generator at each level. Let $P_t = P(\cdot \mid S \le m_t)$ be the level $t$ distribution and $\hat\Lambda_t$ the generator for $-S$ under $P_t$, fitted from the particles at level $t$. The ratio $\theta_{t+1} = P_t(S \le m_{t+1})$ is the length of the interval on which $\hat\Lambda_t \ge -m_{t+1}$, which the fitted generator gives directly, and $\rho = \prod_t \theta_t$ as before. The GBC ladder is the vertical walk with the empirical quantile function at each rung replaced by a fitted one, and the point of the replacement is that the fitted function can be evaluated anywhere, including beyond the range of the sample, with an error governed by the smoothness of the fit rather than the sample size alone (subject to the caution after the Estimation theorem below).

\section{The level set generator}

The deeper use of GBC is in the conditioned draws. A level set sampler produces $X \sim P(\cdot \mid S \le m)$, and this is a conditional law indexed by $m$. By noise outsourcing it is $G(m, \xi)$ for a map $G$ and uniform noise $\xi$, and GBC learns $G$ from pairs $(m, X)$: run the ladder once, record every particle with its level, and fit $G_\phi(m, \xi)$ to generate particles at level $m$ from noise. The fitted $G_\phi$ is then a one shot sampler for every level, with no Markov chain and no mixing time\index{mixing time}, and the ladder can be rerun at any resolution by evaluating $G_\phi$.

Two things must be said about this at once. First, the fitted $G_\phi$ is only as good as the training set, which was produced by the chain, so if the chain failed to mix at level $m$ the generator reproduces the failure: it generates from the component the chain was in. GBC amortises the sampler; it does not improve its coverage. Second, for a discrete space such as permutations the generator must map continuous noise to a discrete structure, which is done by the Gumbel or Plackett Luce parametrisation of permutations, $x = \mathrm{argsort}(\log w_i + \text{Gumbel noise})$ for learned weights $w$, with the weights depending on $m$. This is the standard construction in learning to rank, and it produces a differentiable sampler for permutations whose marginals can be fitted.

What GBC adds, then, is not a new sampler but a representation of the sampler as a function, and functions can be analysed. The map $m \mapsto G_\phi(m, \cdot)$ is the level set sampler as a curve in the space of distributions on permutations, and its regularity in $m$ is a statement about the geometry of the level sets: if the level sets vary continuously with $m$ in the sense that $G_\phi(m, \xi)$ is a continuous function of $m$ for fixed $\xi$, then nearby levels have nearby samplers and the ladder can be coarse. If the map has a discontinuity at some $m^*$, so that a small change of level forces a large change of configuration, there is a phase transition at $m^*$ and the ladder must be fine there. Heuristically, the discontinuities of the level set generator correspond to the phase transitions of the Boltzmann\index{Boltzmann distribution} model, and the GBC representation would make them visible as kinks in a learned function; we do not prove this identification.

\section{Theorems for the GBC ladder}

The following statements are what can be proved about the GBC implementation. They are elementary, but they fix what the method does and does not guarantee.

\begin{theorem}[Representation]
For any finite $\cX$, any $P$ on $\cX$, any $L \ge 0$ and any energy $S$ with levels $m_0 > \dots > m_T$, there exist a nonincreasing function $\Lambda : [0, 1] \to [0, \infty)$ and a measurable $G : \{ m_t \} \times [0, 1] \to \cX$ such that $\Lambda(U)$ has the law of $L(X)$ for $U$ uniform on $[0, 1]$, so that $Z = \int_0^1 \Lambda$, and $G(m_t, \xi) \sim P(\cdot \mid S \le m_t)$ for $\xi$ uniform. Both are unique up to null sets if $G(m_t, \cdot)$ is required to be nondecreasing in a fixed total order on $\cX$.
\end{theorem}

\begin{proof}
$\Lambda$ is the quantile function of $L(X)$ and $G(m_t, \cdot)$ is the quantile function of $P(\cdot \mid S \le m_t)$ in the chosen order. Existence and uniqueness are those of quantile functions.
\end{proof}

\begin{theorem}[Estimation]
Let $\hat\Lambda_M$ be the empirical quantile function of $M$ independent draws of $L(X)$ under $P$. Then for every $s$, $\PP( |\hat\Lambda_M(s) - \Lambda(s)| > \eta ) \le 2 e^{-2 M \delta(s, \eta)^2}$ where $\delta(s, \eta)$ is the probability mass of the interval of likelihood values between $\Lambda(s)$ and $\Lambda(s) \pm \eta$; and the integrated estimate $\int_0^1 \hat\Lambda_M$ equals the sample mean $M^{-1} \sum_i L(X_i)$ exactly.
\end{theorem}

\begin{proof}
The first statement is the Dvoretzky Kiefer Wolfowitz inequality for the empirical distribution function of $L(X)$ transferred to the quantile scale. The second is the identity that the integral of the empirical quantile function is the sample mean.
\end{proof}

The second statement is a caution: integrating the empirical quantile function gives back naive Monte Carlo. The Yakowitz\index{Yakowitz, Krimmel and Szidarovszky} and stratified gains of Chapters~\ref{ch:yakowitz} and~\ref{ch:beyond} come from evaluating $\Lambda$ at chosen points, not from the empirical quantile function of a random sample. A fitted GBC generator, which is a smoothed empirical quantile function, can be evaluated at stratified points, and then the stratified bound applies to the smoothed function with an additional smoothing bias. The gain is real only if the smoothing is accurate, which for a staircase means that the fitted generator locates the jumps, and locating a jump from samples is again a Monte Carlo problem with rate $M^{-1/2}$ in its position. There is no free lunch in the generator either; what it offers is a reusable, differentiable, evaluable representation of what the samples contain.

\begin{theorem}[Amortised ladder]
\label{thm:amortised}
Suppose a level set generator $G_\phi(m, \xi)$ is fitted so that for each level $m_t$ the law of $G_\phi(m_t, \xi)$ is within total variation $\delta$ of $P(\cdot \mid S \le m_t)$. Suppose also that the $T$ ratios equal a common value $\theta$. Then the ladder run with $M$ draws from $G_\phi$ at each rung estimates $\rho$ with relative mean squared error, to first order in $T(1 - \theta)/(\theta M)$ and $T\delta/\theta$, at most $T(1 - \theta)/(\theta M) + O(T^2 \delta^2/\theta^2)$, and costs $T M$ evaluations of $G_\phi$ with no Markov chain steps.
\end{theorem}

\begin{proof}
This is the product estimator\index{product estimator} bound with the sampler's total variation error entering the bias of each ratio additively. A total variation error $\delta$ changes each ratio by at most $\delta$, a relative bias of at most $\delta/\theta$, so the product has relative bias at most $(1 + \delta/\theta)^T - 1 \approx T\delta/\theta$ and squared bias $O(T^2\delta^2/\theta^2)$. For $T$ independent ratio estimates with relative variance $v$ the product has relative variance $(1 + v)^T - 1 = Tv + O(T^2 v^2)$, and $v = (1 - \theta)/(\theta M)$ to first order. Both statements are first order; the exact forms are the two displayed expressions.
\end{proof}

\section{What generators cannot do}

The amortised ladder theorem has $\delta$ in it, and $\delta$ is the error of the generator at each level, which was learned from a chain. If the chain was trapped at level $m$, so that its output was uniform on a component $C$ of the level set rather than on the whole, then the generator learns $C$ and $\delta$ is $1 - |C| / |\{ S \le m \}|$, which is close to one if $C$ is small. No amount of fitting recovers what the training sample did not contain. GBC is a method for turning samples into functions; it is not a method for producing samples that the chain could not produce.

There is one thing a generator can do that a chain cannot, and it is the reason to pursue the GBC implementation. A generator is differentiable in its parameters and in its conditioning variable, and so can be trained with gradient information from the energy. For the queens the energy is a sum of pairwise terms and its expectation under the Plackett Luce law of a weighted argsort is a smooth function of the weights, so one can learn weights $w(m)$ that minimise the expected energy at each level by gradient descent, which is a variational approximation to the level set, and then use the learned $w(m)$ as a proposal for the chain rather than a replacement for it. A chain with a good proposal mixes faster, and the generator, trained on what the chain found, improves the proposal, which improves the chain. This loop, generator proposes and chain corrects, is the GBC analogue of the adaptive Wang Landau\index{Wang Landau algorithm} update, and we propose it as a practical form of vertical counting, to be tested. Its mixing time would still be the mixing time of the corrected chain, and the theorems of Part IV would still apply to it; whether its constants are better is an empirical question we have not tested, since no GBC code accompanies this book.

\section{Monads, functors and the vertical walk}
\label{sec:monads}

The two halves of this book, the Monte Carlo half and the derandomization half, can be placed on one axis by a structural idea that Polson and Sokolov develop in a companion note on the $\lambda$ calculus, monads and generative Bayesian computation \cite{PolsonSokolovMonads}: useful computation is composition, and the constructs that scale are the ones a functor preserves. The idea descends from the same year as the rest of Part~I. Hilbert's Entscheidungsproblem received two answers in 1936, Turing\index{Turing, Alan}'s machine, an imperative device whose primitive notion is the mutation of state, and Church\index{Church, Alonzo}'s $\lambda$ calculus, a functional device whose primitive notion is the application and composition of functions and which has no state at all. Turing proved in 1937 that the two compute the same functions, and the Cook Levin theorem of Chapter~\ref{ch:npcomplete} is indifferent to which is used. But the two traditions handle randomness differently, and the functional handling makes the structure of the vertical walk visible.

A pure function cannot draw a random number. Moggi's \cite{Moggi} device for adding effects to a pure calculus is a monad\index{monad}: a randomized computation returning a value in $A$ is represented as an element of $T A$, the space of probability measures on $A$, and the Giry\index{Giry monad} monad \cite{Giry} supplies the unit $\eta(x) = \delta_x$ and the multiplication $\mu(P)(B) = \int q(B) \, P(dq)$ that marginalises a hierarchical model one level. Its Kleisli category has Markov kernels $A \to T B$ as morphisms and the Chapman Kolmogorov identity as composition, and a probabilistic program is a morphism in it. The relevance to this book is that every object in the Monte Carlo half is such a morphism, and the theorems of the book are statements about what composition does to them.

The level set sampler is a kernel $P_m$ on $\{ S \le m \}$. The vertical walk is the Kleisli composite of the kernels $P_T, P_{T-1}, \dots, P_1$ with the Bernoulli observation at each rung, and the product estimator is the composite's value at the unit. The Lorenz identity of Chapter~\ref{ch:vertical} is the statement that the functor from measures to quantile functions sends the expectation $\EE_P[L]$ to the integral $\int_0^1 \Lambda(s) \, ds$, and the quantile reordering of Chapter~\ref{ch:reorder} is a change of coordinates in the image of that functor. The product estimator variance bound of Chapter~\ref{ch:product} says that the composite of $T$ kernels with relative variance $v$ each has relative variance $(1 + v)^T - 1 = T v + O(T^2 v^2)$; it is the functoriality of the second moment along a Kleisli chain, and the conditional FPRAS theorem of Chapter~\ref{ch:program} is what that functoriality gives when the kernels are polynomial. The amortised ladder theorem of this chapter is the same statement with the kernels replaced by learned generators: a generator $G(m, \xi)$ is a deterministic section of the Kleisli arrow $P_m$, a pure function of base noise $\xi$ whose push forward is the level set law, fitted by the gradient functor of the companion note and evaluated forward at inference time. Generative Bayesian computation is the recognition that the section can be learned once for the whole ladder, and the error $\delta$ in the amortised theorem is the distance between the learned section and the arrow it sections.

The derandomization half has the same form, and seeing it so is the point of the section. A Markov category\index{Markov category} \cite{Fritz} is a monoidal category in which every object can be copied and discarded, and a morphism is deterministic exactly when it commutes with copying: running it twice on one input gives the same answer as copying the output. An $\RP$ algorithm for a language $L$ is a morphism $x \mapsto P_x$ into $T\{0, 1\}$ with $P_x(1) \ge \tfrac12$ when $x \in L$ and $P_x(1) = 0$ otherwise, and it is not deterministic, since a second run can disagree with the first. The question $\Pc = \RP$ asks whether every such morphism that is computable in polynomial time can be replaced by a deterministic polynomial time morphism $x \mapsto \delta_{\Ind\{x \in L\}}$ deciding the same language. In the quantile coordinate of Chapter~\ref{ch:prp} a sufficient condition is the hitting set problem, equivalent to $\mathrm{prRP} = \mathrm{prP}$ (Chapter~\ref{ch:theorems}): a deterministic polynomial time procedure that outputs, for each $x \in L$, a random string in the level set $\{ r : L_x(r) = 1 \}$, which has density at least $\tfrac12$. The two halves are then one question asked of one category. The Monte Carlo half asks whether a given Kleisli arrow, the vertical walk on hard boards, has polynomial cost; the derandomization half asks whether an arrow of polynomial cost can be replaced by one that commutes with copying. Neither question is made easier by the language, and the companion note is explicit that a Markov category does not run faster than a kernel. What the language does is separate the model from its interpreter: the queens board and the energy are the model, a single syntactic object; the swap chain, the block chain, the nested sampler and the learned generator are interpreters of it; and the theorems of Chapter~\ref{ch:theorems} say that every interpreter of that object that runs in polynomial time on every board is a proof that $\NP = \RP$. The model is fixed by Gent, Jefferson and Nightingale. The whole of the program is a search over interpreters.

The functional view also explains an experimental convenience. A pure sampler with a splittable random number generator is reproducible and parallel by construction: for the amortised ladder the level sets can be sampled on separate cores with independent streams, the ladder estimates combined by a map and a reduce, and the result does not depend on the scheduling. The consensus Monte Carlo of Scott and coauthors \cite{ScottConsensus} is the same map and reduce over data shards. For the amortised ladder of this chapter, whose rungs draw independently from $G_\phi$, the map is over rungs and the reduce is the product, and the error analysis of Chapter~\ref{ch:product} is the statement that the reduce amplifies the relative variance of the map by a factor of about $T$, to first order. This is the practical content of functoriality, and it is why the experiments of Part~IV, which were run on a single machine, could be distributed across many machines as independent runs, though the particle walk itself is sequential across rungs.
